\section{Auswertung}
\label{sec:Auswertung}
\subsection{Magnetfeldstärke an der Probe}
\label{sec:1}
Da für die Bestimmung der effektiven Masse und somit für die Auswertung der 
Messwerte die Kenntnis der Magnetfeldverteilung entlang der Probe notwendig ist,
wurde diese zunächst mit einer Hall-Sonde vermessen. Die Messwerte wurden in
\autoref{fig:magnetfeldverteilung} dargestellt und mit einem quadratischen Fit
versehen, um die Verteilung zu beschreiben. Die gefundenen Koeffizienten des Fits
sind in der Legende der Grafik angegeben. Es ist zu erkennen, dass die Messdaten
gut durch den quadratischen Fit beschrieben werden, jedoch wurden die Werte bei
jeweils $\qty{80}{mm}$ sowie $\qty{120}{mm}$ nicht berücksichtigt, da diese 
offensichtlich Ausreißer darstellen, da das Magnetfeld dort nahezu Null ist. Die
Messwerte sind \autoref{tab:magnetfeldverteilung} zu entnehmen.

\begin{figure}
    \centering
    \includegraphics[width=0.8\textwidth]{build/magnetfeldverteilung.png}
    \caption{Messwerte der Magnetfeldverteilung entlang der Probe mit quadratischem Fit.}
    \label{fig:magnetfeldverteilung}
\end{figure}

\begin{table}[H]
    \centering
    \caption{Messwerte der Magnetfeldverteilung entlang der Probe.}
    \label{tab:magnetfeldverteilung}
    \begin{tabular}{c|c}
        Position $x$ (mm) & Magnetfeld $B$ (mT) \\
        \hline
        100 & 422\\
        95 & 402\\
        90 & 300\\
        85 & 104\\
        80 & 24\\
        105 & 380\\ 
        110 & 223\\
        115 & 56\\
        120 & 13\\
    \end{tabular}
\end{table}
\noindent Der Hochpunkt der Ausgleichsrechnung und damit die maximale Magnetfeldstärke ergibt sich zu: 
\begin{equation}
    B_\text{max} = \qty{424+-9}{mT}
\end{equation}


\subsection{Bestimmung der effektiven Masse}
Die Bestimmung der effektiven Masse erfolgt durch die Auswertung der gemessenen
Faraday-Rotationswinkel der freien Ladungsträger in Abhängigkeit von der
Wellenlänge. Zuallererst werden alle gemessenen Winkel durch die Dicke der jeweiligen
Probe geteilt und somit skaliert. Dann werden die skalierten Rotationswinkel der
undotierten Probe von den beiden gemesenen Rotationswinkeln bei den
$1.2\cdot10^{18} \text{cm}^{-3}$ und $2.8\cdot10^{18} \text{cm}^{-3}$ dotierten
Proben subtrahiert, um die freien Faraday-Rotationswinkel zu erhalten. Anschließend
wurden diese gegen das Quadrat der Wellenlänge aufgetragen und mit einem linearen
Fit versehen, um die Steigung zu bestimmen. Die effektive Masse kann dann aus der
Steigung des Fits berechnet werden.  
\vspace{0.5cm}
\\
\noindent Die gemessenen Faraday-Winkel sind \autoref{tab:rotationswinkel1} und
\autoref{tab:rotationswinkel2} zu entnehmen. Die Rotationswinkel der zwei Proben
sowie der undotierten Probe, sind die halbierten Differenzen der aufgenommenen 
Winkel bei der jeweiligen Wellenlänge.
\begin{table}[H]
    \centering
    \caption{Spezifikationen der Proben}
    \label{tab:proben}
    \begin{tabular}{c|c|c}
        Probe & Dotierung & Dicke (mm) \\
        \hline
        1 & $1.2\cdot10^{18} \text{cm}^{-3}$ & 1.36 \\
        2 & $2.8\cdot10^{18} \text{cm}^{-3}$ & 1.29 \\
        3 & undotiert & 5.11 \\
    \end{tabular}
\end{table}
\begin{table}[H]
    \centering
    \caption{Gemessene Rotationswinkel sowie freie Rotation für Probe 1 (1.2$\cdot$10$^{18}$ cm$^{-3}$ Dotierung).}
    \label{tab:rotationswinkel1}
    \begin{tabular}{c|c|c|c}
        Wellenlänge $\lambda$ (nm) & $\theta_1$ / rad & $\theta_2$ / rad & $\theta_\text{frei}$ / rad \\
        \hline
        1,06 &  8,17e-2 & -6,20e-2 &   72,2 \\
        1,29 &  5,21e-2 &  1,23e-1 &   14,3 \\
        1,45 &  2,01e-2 &  9,96e-2 &   -4,7 \\
        1,72 &  2,36e-2 &  9,49e-1 & -168,0 \\ % Ausreißer Filterwechsel
        1,96 &  5,92e-2 & -8,46e-1 &  209,0 \\ % Ausreißer Filterwechsel
        2,16 &  6,20e-2 &  4,61e-2 &   36,5 \\
        2,34 &  6,18e-2 &  3,39e-2 &   38,8 \\
        2,51 &  6,60e-2 &  2,75e-2 &   43,2 \\
        2,65 &  5,44e-2 &  4,29e-2 &   31,6 \\
    \end{tabular}
\end{table}
\begin{table}[H]
    \centering
    \caption{Gemessene Rotationswinkel sowie freie Rotation für Probe 2 (2.8$\cdot$10$^{18}$ cm$^{-3}$ Dotierung).}
    \label{tab:rotationswinkel2}
    \begin{tabular}{c|c|c|c}
        Wellenlänge $\lambda$ (nm) & $\theta_1$ / rad & $\theta_2$ / rad & $\theta_\text{frei}$ / rad \\
        \hline
        1.06     & 9.85e-02 & -6.20e-02 & 8.85e+01\\
        1.29     & 7.71e-02 & 1.23e-01 & 3.58e+01\\
        1.45     & 7.77e-02 & 9.96e-02 & 4.07e+01\\
        1.72     & 8.49e-02 & 9.49e-01 & -1.20e+02\\
        1.96     & 9.96e-02 & -8.46e-01 & 2.43e+02\\
        2.16     & 1.02e-01 & 4.61e-02 & 7.01e+01\\
        2.34     & 7.36e-02 & 3.39e-02 & 5.04e+01\\
        2.51     & 1.56e-01 & 2.75e-02 & 1.15e+02   \\
        2.65     & 1.18e-01 & 4.29e-02 & 8.29e+01   \\
    \end{tabular}   
\end{table}
Für die Berechnung der effektiven Masse werden die folgenden Konstanten benötigt:
\begin{align*}
    \label{eqn:1}
    e_0      &= 1{,}602 \times 10^{-19} \, \text{C} && \text{(Elementarladung)} \\
    \epsilon_0 &= 8{,}854 \times 10^{-12} \, \text{F/m} && \text{(Elektrische Feldkonstante)} \\
    c        &= 2{,}998 \times 10^{8} \, \text{m/s} && \text{(Lichtgeschwindigkeit)} \\
    n        &= 3{,}3 && \text{(Brechungsindex von GaAs)} \\
    m_e      &= 9{,}109 \times 10^{-31} \, \text{kg} && \text{(Ruhemasse des freien Elektrons)}
\end{align*}
\noindent Nun werden die freien Rotationswinkel gegen das Quadrat der Wellenlänge
aufgetragen und mit einem linearen Fit versehen, um die Steigung zu bestimmen.
In der Gleichung 
\begin{equation*}
    \theta_\text{frei} = \frac{e^3}{8\pi^2\epsilon_0 c^3}\frac{N B_\text{max}}{n m^{*2}}\lambda^2
\end{equation*}
kann der Term 
\begin{equation*}
    a = \frac{e^3}{8\pi^2\epsilon_0 c^3}\frac{N B_\text{max}}{n m^{*2}}
\end{equation*}
als Steigung der Ausgleichsgeraden identifiziert werden. Durch Umstellen und 
Einsetzen der Konstanten aus \autoref{eqn:1} sowie der in \autoref{sec:1} berechneten 
maximale Magnetfeldstärke am Ort der Probe, lässt scih auf die effektive Masse schließen.
\begin{equation}
    \label{eqn:2}
    m^* = \sqrt{\frac{e^3}{8\pi^2\epsilon_0 c^3}\frac{N B_\text{max}}{n a}}
\end{equation}
\begin{figure}[H]
    \centering
    \caption{Freie Rotationswinkel gegen das Quadrat der Wellenlänge mit linearem Fit.}
    \includegraphics[width=0.8\textwidth]{build/faraday_rotation.png}
    \label{fig:faradayrotation}
\end{figure} 
\noindent Bei der ersten Probe wurde der erste Messwert und bei den beiden anderen
Proben der dritte und vierte Messwert nicht in die Regression mit einbezogen, da diese 
offensichtlich Ausreißer darstellen.
Die Parameter der Ausgleichsrechnung ergeben sich zu
\begin{align*}
    a_{1.2dot} &= \qty{6.8+-2.4e+12}{rad/m^2} \\
    a_{2.8dot} &= \qty{6+-5e+12}{rad/m^2}
\end{align*}
Durch Einsetzen in \autoref{eqn:2} ergeben sichdie effektiven Massen als
\begin{align}
    m^*_{1.2dot} &= \qty{7.0+-1.2e-32}{kg}, \\
    m^*_{2.8dot} &= \qty{1.2+-0.5e-31}{kg}.
\end{align}